\section{Choice of a Cascade-Graph Model}
  \label{sec:model}

  \subsection{Requirements on the cascade graphs}
    \label{sec:model:requirements}

    The following properties are necessary for a family $\{G(n)\}_{n \geq 1}$ to
    be consistent with the projective cascade picture.

    \medskip
    \noindent\textbf{(R1) Vertex-transitivity.}
    The projection must not privilege any node; $G(n)$ must be vertex-transitive
    or quasi-homogeneous, so that $\lambda_2(n)$ is a well-defined global quantity.

    \medskip
    \noindent\textbf{(R2) Bounded degree.}
    The local relational flux capacity is finite, so $d(n)$ must remain bounded
    (or grow at most logarithmically) along the cascade.

    \medskip
    \noindent\textbf{(R3) Non-vanishing algebraic connectivity.}
    $\lambda_2(n) > 0$ must hold for all $n$, ensuring a coherent global
    projection.

    \medskip
    \noindent\textbf{(R4) Controlled spectral gap.}
    $\lambda_2(n)$ must admit explicit analytic lower bounds.

    \medskip
    \noindent\textbf{(R5) Constructibility as an explicit sequence.}
    The family must admit an explicit indexing compatible with the cascade rank $n$.

  \subsection{Candidate families and selection of LPS graphs}
    \label{sec:model:candidates}

    Random regular graphs satisfy (R1)--(R4) only in probability and not
    constructively, failing (R5).
    Cayley graphs of finite groups satisfy (R1) and (R4) for families with
    controlled character sums, but generic sequences need not satisfy (R3).
    Explicit expander families satisfy all five requirements by construction.
    Among expanders, \emph{near-Ramanujan} families are optimal: the Alon--Boppana
    theorem~\cite{Alon1986} bounds $\lambda_2 \leq 1 - 2\sqrt{d-1}/d + o(1)$, and
    Ramanujan graphs achieve this bound exactly.

    The Lubotzky--Phillips--Sarnak (LPS) construction~\cite{LPS1988} provides the
    only currently known explicit infinite families of Ramanujan graphs, and is
    the natural canonical choice satisfying all five requirements.

  \subsection{LPS graphs}
    \label{sec:model:lps}

    Let $p$ and $q$ be distinct odd primes, both congruent to $1 \pmod{4}$.
    The LPS graph $X^{p,q}$ is a $(p+1)$-regular Cayley graph on
    $\mathrm{PSL}(2, \mathbb{F}_q)$, with
    \begin{equation}
      \label{eq:lps-order}
      \bigl|V(X^{p,q})\bigr|
      \;=\; \tfrac{1}{2}\,q(q^2 - 1).
    \end{equation}
    The Ramanujan property gives, for the normalised Laplacian,
    \begin{equation}
      \label{eq:ramanujan-laplacian}
      \lambda_2\bigl(X^{p,q}\bigr)
      \;\geq\; 1 - \frac{2\sqrt{p}}{p+1},
    \end{equation}
    uniformly in $q$.
    All five requirements (R1)--(R5) are satisfied.

    \begin{remark}[Laplacian normalisation]
  Throughout this paper, $L_G = I - A/d$ denotes the normalised graph
  Laplacian, with eigenvalues in $[0,2]$.
  The Cheeger inequality~\eqref{eq:cheeger-ineq} and the Ramanujan
  bound~\eqref{eq:ramanujan-laplacian} use this convention consistently.
    \end{remark}

  \subsection{Identification of the cascade index}
    \label{sec:model:index}

    The physical identification is
    \begin{equation}
      \label{eq:index-id}
      n \;\sim\; \bigl|V\bigl(G(n)\bigr)\bigr|,
    \end{equation}
    interpreting the cascade depth as proportional to the number of active
    relational nodes.
    For the LPS family at fixed $p$,
    \begin{equation}
      \label{eq:index-lps}
      n \;\sim\; \tfrac{1}{2}\,q(q^2 - 1) \;\sim\; q^3/2
      \qquad (q \to \infty),
    \end{equation}
    so $q \sim (2n)^{1/3}$.
    Alternative identifications ($n \sim |E(G(n))|$, $n \sim \log|V|$) are
    discussed in Section~\ref{sec:discussion}.
